\section{Discussion}
The physical result is a hierarchy of structure rather than a single phase map.  The two-dimensional exact data show a narrow uniform-SC strip, an extended finite-$q$ region, and a curved Normal--superconducting envelope; the spectral analysis then resolves the FFLO sector into gapped conventional and topological bulk supports.  Varying chemical potential reorganizes the coordinate support of those states and can remove trusted tFFLO observations from an individual sampled slice.  Topology therefore supplies information that a thermodynamic FFLO label alone cannot encode.

Methodologically, reliability-aware active learning differs from a black-box phase classifier because physical labels are obtained from the BdG calculation at every layer.  Numerical reliability is stored separately from phase identity, topology is masked outside reliable gapped FFLO states, and uncertainty affects where the next calculation is performed rather than what phase is declared.  This organization makes the expensive minimization and spectral tests selectively deployable without turning the surrogate into a physics solver.

The location-count economy does not establish wall-time acceleration, because no matched exhaustive timing baseline is archived, and it does not establish boundary enrichment, because the candidate-wide boundary denominator is unavailable.  The stopping settings diagnose a particular numerical workflow rather than universal physical scales.  Likewise, the three-dimensional evidence is discrete and nonuniformly sampled, so it supports higher-dimensional exploration but neither phase volumes nor a closed continuous atlas.

The broader significance is a route for coupling direct BdG numerical objectives to adaptive exploration while preserving applicability masks and failure states.  Exploratory learned-acquisition studies have not established superiority over the production acquisition.  A future physical world model would need spectral descriptors, calibrated uncertainty, and acquisition utility tied to physical validity; these ingredients remain an outlook rather than a validated result here.

\paragraph{Data availability.} The numerical data and source configurations are preserved in the project archive. No public repository or persistent identifier has been established for this review version.
